\section{Memory Issues}

\begin{frame}
  \frametitle{Usual Memory Issues}
  \begin{itemize}
    \item Programming (almost) always involves accessing memory
    \item If done incorrectly, a large variety of errors can be triggered
    \begin{itemize}
      \item Segmentation Faults can happen when accessing invalid memory
            addresses (NULL pointers or use-after-free for instance)
      \item Buffer Overflows can happen if accessing a buffer outside its
            boundaries
      \item Memory Leaks when allocating memory and forgetting to free it after
            usage
    \end{itemize}
    \item Fortunately, there are tools to debug these errors
  \end{itemize}
\end{frame}

\begin{frame}[fragile]
  \frametitle{Segmentation Faults}
  \begin{itemize}
    \item Segmentation Faults are generated by the kernel when a program tries to
          access a memory area that it is not allowed to or to access it in an
          incorrect way
    \begin{itemize}
      \item Might be generated by a write on a read only memory zone
      \item Can also be triggered when trying to execute memory that is not
            executable
    \end{itemize}
  \end{itemize}

  \begin{block}{}
    \begin{minted}[fontsize=\small]{c}
int *ptr = NULL;
*p = 1;
    \end{minted}
  \end{block}

  \begin{itemize}
    \item Execution will yield a \code{Segmentation fault} message in the
          terminal
  \end{itemize}
  \vspace{0.2cm}

  \begin{block}{}
    \begin{minted}[fontsize=\small]{console}
$ ./program
Segmentation fault
    \end{minted}
  \end{block}
\end{frame}

\begin{frame}[fragile]
  \frametitle{Buffer Overflows}
  \begin{itemize}
    \item Buffer Overflows are easily triggered when accessing an array outside
          of its boundaries (most often past the end)
    \item Such access might generate a crash or not depending on the access
    \begin{itemize}
      \item Writing past the end of a \code{malloc()}'ed array will most often
            overwrite the malloc data structure leading to corruption
      \item Writing past the end of an array allocated on the stack can corrupt
            data on the stack
      \item Reading past the end of an array might generate a segfault but not
            always, this depends on the area of memory that is accessed
    \end{itemize}
  \end{itemize}
  \vspace{0.2cm}
  \begin{block}{}
    \begin{minted}[fontsize=\small]{c}
uint32_t *array = malloc(10 * sizeof(*array));
array[10] = 0xDEADBEEF;
    \end{minted}
  \end{block}
\end{frame}

\begin{frame}[fragile]
  \frametitle{Memory Leaks}
  \begin{itemize}
    \item Memory leaks are another class of memory errors that will not directly
          trigger a crash but will exhaust the system memory (sooner or later)
    \item This happens when allocating memory in your program and not releasing
          it after using it
    \item Can trigger in production when the program runs for a very long time
    \begin{itemize}
      \item Better to debug that kind of problem early in the development
            process
    \end{itemize}
  \end{itemize}
  \vspace{0.2cm}
  \begin{block}{}
    \begin{minted}[fontsize=\small]{c}
void func1(void) {
  uint32_t *array = malloc(10 * sizeof(*array));
  do_something_with_array(array);
}
    \end{minted}
  \end{block}
\end{frame}

\begin{frame}
  \frametitle{Valgrind (1/2)}
  \begin{columns}[T]
    \column{0.8\textwidth}
    \begin{itemize}
      \item \href{https://valgrind.org/}{Valgrind} is an instrumentation
            framework for building dynamic analysis tools
      \item \code{valgrind} is also a tool that is based on this framework
            and provides a memory error detector, heap profilers and others
            profilers.
      \item It is supported on all the popular platforms: Linux on x86, x86\_64,
            arm (armv7 only), arm64, mips32, s390, ppc32 and ppc64.
    \end{itemize}
    \column{0.2\textwidth}
    \includegraphics[width=\textwidth]{../common/valgrind1.png}
  \end{columns}
\end{frame}

\begin{frame}[fragile]
  \frametitle{Valgrind (2/2)}
  \begin{columns}[T]
    \column{0.8\textwidth}
    \begin{itemize}
      \item Works by adding its own instrumentation to your code and
            then running in on its own virtual cpu core. Significantly slows
            down execution, and thus is suited for debugging and profiling
      \item {\em Memcheck} is the defaut {\em valgrind} tool and it detects
            memory-management errors
      \begin{itemize}
        \item Access to invalid memory zones, use of uninitialized values,
              memory leaks, bad freeing of heap blocks, etc
        \item Can be run on any application, no need to recompile them
      \end{itemize}
    \end{itemize}
      \begin{block}{}
        \begin{minted}[fontsize=\small]{console}
    $ valgrind --tool=memcheck --leak-check=full <program>
        \end{minted}
      \end{block}
    \column{0.2\textwidth}
    \includegraphics[width=\textwidth]{../common/valgrind2.png}
  \end{columns}
\end{frame}

\begin{frame}[fragile]
  \frametitle{Valgrind Memcheck usage and report}
  \begin{block}{}
    \begin{minted}[fontsize=\tiny]{console}
$ valgrind --leak-check=full ./mem_leak
==202104== Memcheck, a memory error detector
==202104== Copyright (C) 2002-2017, and GNU GPL'd, by Julian Seward et al.
==202104== Using Valgrind-3.18.1 and LibVEX; rerun with -h for copyright info
==202104== Command: ./a.out
==202104==
==202104== Conditional jump or move depends on uninitialised value(s)
==202104==    at 0x109161: do_actual_jump (in /home/user/mem_leak)
==202104==    by 0x109187: compute_address (in /home/user/mem_leak)
==202104==    by 0x1091A2: do_jump (in /home/user/mem_leak)
==202104==    by 0x1091D7: main (in /home/user/mem_leak)
==202104==
==202104== HEAP SUMMARY:
==202104==     in use at exit: 120 bytes in 1 blocks
==202104==   total heap usage: 1 allocs, 0 frees, 120 bytes allocated
==202104==
==202104== LEAK SUMMARY:
==202104==    definitely lost: 120 bytes in 1 blocks
==202104==    indirectly lost: 0 bytes in 0 blocks
==202104==      possibly lost: 0 bytes in 0 blocks
==202104==    still reachable: 0 bytes in 0 blocks
==202104==         suppressed: 0 bytes in 0 blocks
==202104== Rerun with --leak-check=full to see details of leaked memory
    \end{minted}
  \end{block}
\end{frame}


\begin{frame}[fragile]
  \frametitle{Using GDB with Memcheck}
  \begin{itemize}
    \item {\em valgrind} allows to attach with GDB to the process that is
          currently analyzed.
  \end{itemize}
  \begin{block}{}
    \begin{minted}[fontsize=\tiny]{console}
$ valgrind --tool=memcheck  --leak-check=full --vgdb=yes --vgdb-error=0 ./mem_leak
    \end{minted}
  \end{block}

  \begin{itemize}
    \item Then attach gdb to the valgrind gdbserver using {\em vgdb}
  \end{itemize}
  \begin{block}{}
    \begin{minted}[fontsize=\tiny]{console}
$ gdb ./mem_leak
(gdb) target remote | vgdb
    \end{minted}
  \end{block}

  \begin{itemize}
    \item If valgrind detects an error, it will stop the execution and break
          into GDB.
  \end{itemize}

  \begin{block}{}
    \begin{minted}[fontsize=\tiny]{console}
(gdb) continue
Continuing.

Program received signal SIGTRAP, Trace/breakpoint trap.
0x0000000000109161 in do_actual_jump (p=0x4a52040) at mem_leak.c:5
5   if (p[1])
(gdb) bt
#0  0x0000000000109161 in do_actual_jump (p=0x4a52040) at mem_leak.c:5
#1  0x0000000000109188 in compute_address (p=0x4a52040) at mem_leak.c:11
#2  0x00000000001091a3 in do_jump (p=0x4a52040) at mem_leak.c:16
#3  0x00000000001091d8 in main () at mem_leak.c:27
    \end{minted}
  \end{block}
\end{frame}

\begin{frame}[fragile]
  \frametitle{libefence (1/2)}
  \begin{itemize}
    \item {\em libefence} is more lightweight than {\em valgrind} but less precise
    \item Allows to catch two types of common memory errors
    \begin{itemize}
      \item Buffer overflows and use after free
    \end{itemize}
    \item {\em libefence} will actually trigger a segfault upon the first error
          encountered in order to generate a coredump.
    \item Uses a shared library that can either be linked with statically
          (\code{-lefence}) or preloaded using \code{LD_PRELOAD}.
  \end{itemize}

  \begin{block}{}
    \begin{minted}[fontsize=\small]{console}
$ gcc -g program.c -o program
$ LD_PRELOAD=libefence.so.0.0 ./program

  Electric Fence 2.2 Copyright (C) 1987-1999 Bruce Perens <bruce@perens.com>
Segmentation fault (core dumped)
    \end{minted}
  \end{block}

\end{frame}

\begin{frame}[fragile]
  \frametitle{libefence (2/2)}
  \begin{itemize}
    \item Upon segfault, a coredump will be generated in the current directory
    \item This coredump can be opened with GDB and will pinpoint the exact
          location where the error did happened
  \end{itemize}
  \begin{block}{}
    \begin{minted}[fontsize=\small]{console}
$ gdb ./program core-program-3485
Reading symbols from ./libefence...
[New LWP 57462]
[Thread debugging using libthread_db enabled]
Using host libthread_db library "/lib/x86_64-linux-gnu/libthread_db.so.1".
Core was generated by `./libefence'.
Program terminated with signal SIGSEGV, Segmentation fault.
#0  main () at libefence.c:8
8    data[99] = 1;
(gdb) 
    \end{minted}
  \end{block}
\end{frame}

\setuplabframe
{Debugging Memory Issues}
{
  Debug various memory issues using specific tooling
  \begin{itemize}
    \item Memory leak and misbehavior detection with {\em valgrind} and
            {\em vgdb}.
  \end{itemize}
}